\documentclass[letterpaper,11pt]{article}
\usepackage[top=0.8in, bottom=0.8in, left=0.9in, right=0.9in]{geometry}
\usepackage[hidelinks]{hyperref}
\usepackage{lmodern}
\usepackage{parskip}
\usepackage{microtype}
\setlength{\parskip}{6pt}
\setlength{\parindent}{0pt}
\begin{document}
\begin{flushleft}
\textbf{\large Ajay Venkatesh} \\
Brooklyn, NY \\
+1 516 585 9013 \\
\href{mailto:av3855@nyu.edu}{av3855@nyu.edu}
\end{flushleft}
\vspace{10pt}
May 21, 2026
\vspace{6pt}

Hiring Manager \\
The University of Chicago Booth School of Business

\vspace{10pt}

Dear Hiring Manager,

My name is Ajay Venkatesh. I'm finishing my M.S. in Computer Engineering at NYU, with production software experience at PwC in Bengaluru, a summer SDE internship at Amazon, and ongoing on-campus work at NYU's Novel AI Technologies. I'm writing about the Research Professional role at the University of Chicago Booth School of Business.

The role's combination of data collection, statistical analysis, and writing statistical software in support of faculty research lines up with my engineering work and academic background. My graduate coursework spans \textbf{Machine Learning}, \textbf{Deep Learning}, \textbf{Big Data}, and \textbf{Advanced Java}, with a Big Data project that built a \textbf{PySpark / HDFS} pipeline over millions of flight records, training Random Forest and Gradient Boosted Tree models with rigorous EDA, multivariate analysis, and feature engineering.

On research support and reproducible engineering: I partnered with researchers at NYU's Novel AI Technologies on an NSF-funded health platform, engineering reproducible \textbf{ETL pipelines} transforming \textbf{NoSQL} into \textbf{PostgreSQL} with concurrency control. The same discipline, careful data infrastructure that supports ongoing research questions, is what faculty research depends on.

On software for analysis: at Amazon I shipped backend services in \textbf{Java} with low-level design patterns and rigorous testing through \textbf{Mockito}. I'm comfortable writing clean, well-documented code that other researchers can read, modify, and build on, the kind of statistical software that compounds across studies.

Stack-wise: \textbf{Python}, \textbf{SQL}, \textbf{Pandas}, \textbf{NumPy}, \textbf{scikit-learn}, \textbf{PySpark}, plus \textbf{PostgreSQL} and \textbf{MySQL} for data work. Comfortable with statistical methods and at home with the careful, reproducible discipline that research engineering requires.

What pulls me toward Chicago Booth is the influence of the research itself. Business and economic research at this level shapes how policy, finance, and management are actually understood, and contributing to that pipeline at the data-and-software layer is meaningful work beyond any single product.

I'd welcome the chance to discuss further. Thanks for considering my application.

\vspace{12pt}
Sincerely, \\
Ajay Venkatesh

\end{document}